\documentclass[11pt]{article}
\usepackage[margin=1in]{geometry}
\usepackage{amsmath, amssymb}
\usepackage{xcolor}
\usepackage{enumitem}
\usepackage{titlesec}
\usepackage[colorlinks=true, linkcolor=blue!50!black, urlcolor=blue!50!black, citecolor=blue!50!black]{hyperref}
\usepackage{fancyhdr}

\definecolor{codeblue}{HTML}{213A87}
\definecolor{accentred}{HTML}{BF1F1F}

\titleformat{\section}{\Large\bfseries\color{codeblue}}{\thesection}{0.8em}{}
\titleformat{\subsection}{\large\bfseries\color{codeblue}}{\thesubsection}{0.6em}{}

\pagestyle{fancy}
\fancyhf{}
\lhead{\small Shortlisted Bibliography --- PAT-scan Literature}
\rhead{\small \thepage}
\renewcommand{\headrulewidth}{0.4pt}

\setlength{\parindent}{0pt}
\setlength{\parskip}{4pt}

\title{Shortlisted Bibliography \\[4pt]
  \large PAT-scan Literature --- User-Curated Stack (Stack 3)}
\author{Claude Opus 4.7}
\date{May 31, 2026}

\newcommand{\bibentry}[6]{%
  \noindent\textbf{[#1]}~#2.~\textit{#3}.~#4, #5.~%
  \href{https://doi.org/#6}{\texttt{doi:#6}}\par\vspace{6pt}
}

\begin{document}
\maketitle
\thispagestyle{fancy}

\begin{center}
\small\itshape
Eight papers enriched via OpenAlex from a user-curated tree of three categories: \\
Mathematical Foundations, Statics-based Methods, Vibrations-based Methods.
\end{center}

\vspace{1em}
\hrule
\vspace{1em}

%% ============================================================
\section*{Summary}

\begin{tabular}{@{}lrl@{}}
\textbf{Category} & \textbf{Count} & \textbf{Year range} \\
\hline
Mathematical Foundations & 1 & 1994 \\
Statics-based Methods    & 5 & 2003--2022 \\
Vibrations-based Methods & 2 & 2007--2022 \\
\hline
\textbf{Total}           & \textbf{8} & \textbf{1994--2022} \\
\end{tabular}

\vspace{1em}

%% ============================================================
\section{Mathematical Foundations}

\bibentry{1}{Raghavan, K.\ and Yagle, A.\,E.}{Forward and inverse problems in elasticity imaging of soft tissues}{IEEE Transactions on Nuclear Science}{1994}{10.1109/23.322961}

\textit{University of Michigan--Ann Arbor.} Cited by 99. \quad
Foundational formulation casting the elasticity inverse problem as a linear system via finite differences, with adaptive piston deformations. Establishes the consistency between forward and inverse problems exploited by all later work in this stack.

%% ============================================================
\section{Statics-based Methods}

\bibentry{2}{Konofagou, E.\,E.\ and Harrigan, T.\,P.}{Palpation tomography --- a new technique for modulus estimation in elastography}{IEEE Ultrasonics Symposium Proceedings}{2003/2004}{10.1109/ultsym.2003.1293486}

\textit{Columbia University; Exponent (United States).} \quad
Uses multiple smaller quasi-static load cases (inspired by clinical palpation) instead of one whole-boundary load, increasing the measurement-to-parameter ratio and reducing noise sensitivity in modulus reconstruction.

\bibentry{3}{Liu, Y., Sun, L., and Wang, G.}{Tomography-based 3-D anisotropic elastography using boundary measurements}{IEEE Transactions on Medical Imaging}{2005}{10.1109/tmi.2005.857232}

\textit{University of Iowa; University of Iowa Hospitals and Clinics.} Cited by 41. \quad
Anisotropic elastography framework using static displacements on the entire boundary plus partial surface force data, with adjoint-method gradients accelerating the optimization-based reconstruction.

\bibentry{4}{Olson, L.\,G.\ and Throne, R.\,D.}{An inverse problem approach to stiffness mapping for early detection of breast cancer}{Inverse Problems in Science and Engineering}{2012}{10.1080/17415977.2012.700710}

\textit{Rose--Hulman Institute of Technology.} Cited by 13. \quad
3D finite-element reformulation of the boundary force/deflection inverse problem, with a genetic algorithm identifying interior tissue moduli. Detects 1\,cm tumours at 23\,dB SNR.

\bibentry{5}{Mei, Y., Wang, S., Shen, X., Rabke, S., and Goenezen, S.}{Mechanics based tomography: a preliminary feasibility study}{Sensors (MDPI)}{2017}{10.3390/s17051075}

\textit{Texas A\&M University.} Cited by 19.\quad Open access (gold). \quad
Coins the term \emph{Mechanics Based Tomography} (MBT). Uses force sensors plus DIC-measured boundary displacements --- no interior measurements --- to reconstruct shear-modulus maps of stiff inclusions in a soft background.

\bibentry{6}{Zhang, E., Dao, M., Karniadakis, G.\,E., and Suresh, S.}{Analyses of internal structures and defects in materials using physics-informed neural networks}{Science Advances}{2022}{10.1126/sciadv.abk0644}

\textit{Brown University; MIT; Nanyang Technological University.} Cited by 344.\quad Open access (gold). \quad
PINN framework for identifying unknown geometry and material parameters with a differentiable, trainable parameterization. Validated across linear elasticity, hyperelasticity, and plasticity --- recovers void/inclusion shape, size, location, and modulus.

%% ============================================================
\section{Vibrations-based Methods}

\bibentry{7}{Peters, A.}{Digital image elasto-tomography: mechanical property reconstruction from surface measured displacement data}{Ph.D.\ thesis, University of Canterbury}{2007}{10.26021/2670}

\textit{University of Canterbury.}\quad Open access (green). \quad
Introduces Digital Image Elasto-Tomography (DIET): steady-state surface motion captured by a camera array drives a finite-element-based stiffness reconstruction, combining deterministic and stochastic algorithms. Targets breast-cancer screening as a low-cost, high-contrast alternative to mammography.

\bibentry{8}{Feng, B.\,T., Ogren, A.\,C., Daraio, C., and Bouman, K.\,L.}{Visual vibration tomography: estimating interior material properties from monocular video}{IEEE/CVF Conference on Computer Vision and Pattern Recognition (CVPR)}{2022}{10.1109/cvpr52688.2022.01575}

\textit{California Institute of Technology.} Cited by 12.\quad Open access (green). \quad
Estimates spatially-varying Young's modulus and density throughout a 3D object from a single monocular video, by extracting image-space modes from sub-pixel surface vibration and inverting to material properties. Validated on simulated and real videos.

\vspace{2em}
\hrule
\vspace{0.5em}

{\footnotesize\itshape
Generated by the Stack 3 cap-stone skill (\texttt{/identify-papers-user}) over the curated tree at \\
\texttt{paper-stack-three-test-two/shortlisted/}. Source: OpenAlex (resolution mode: search) on May 31, 2026.
}

\end{document}
